% Section declaración - DATA SCIENCE / ANÁLISIS DE DATOS
\par{

	Economista con fuerte orientación hacia la \textbf{ciencia de datos} y el análisis cuantitativo aplicado. Especializado en procesamiento, depuración y análisis de bases de datos complejas utilizando \textbf{R}, \textbf{Python} y \textbf{SQL}, con capacidad de transformar datos en insights accionables para la toma de decisiones.

	Experiencia en desarrollo de dashboards interactivos, elaboración de reportes automatizados, análisis estadístico avanzado y visualización de información cuantitativa. Dominio de técnicas de Machine Learning aplicadas a economía, proyecciones econométricas y análisis de series temporales. Familiarizado con metodologías ágiles y control de versiones (Git).

	Busco oportunidades en equipos de Data Science, Business Intelligence o Analytics donde pueda aplicar mi formación económica y mis habilidades técnicas para desarrollar soluciones basadas en datos, contribuir con modelamiento predictivo, automatización de procesos analíticos y generación de valor a través del análisis de información.

	%%%%%%%%%%%%%%%%%%%%%%%%%%%%%%%%%%%%%%%%
	% VERSIÓN ALTERNATIVA - Business Intelligence
	%%%%%%%%%%%%%%%%%%%%%%%%%%%%%%%%%%%%%%%%
	%Economista especializado en Business Intelligence y análisis de datos con experiencia en desarrollo de dashboards, reportería automatizada y análisis de KPIs. Experto en transformar datos complejos en visualizaciones intuitivas que faciliten la toma de decisiones estratégicas.

	%Dominio de herramientas de visualización (Power BI, Tableau, R Shiny) y lenguajes de programación orientados a datos (R, Python, SQL). Capacidad de diseñar e implementar pipelines de datos, crear reportes ejecutivos y comunicar insights técnicos a audiencias no técnicas de manera efectiva.

	%Busco contribuir en áreas de Business Intelligence, Analytics o Reporting donde pueda aplicar mi visión analítica para optimizar procesos de negocio, desarrollar tableros de control y generar inteligencia de negocio que impulse el crecimiento organizacional.

	%%%%%%%%%%%%%%%%%%%%%%%%%%%%%%%%%%%%%%%%
	% VERSIÓN ALTERNATIVA - Data Analyst entry-level
	%%%%%%%%%%%%%%%%%%%%%%%%%%%%%%%%%%%%%%%%
	%Recién egresado de Economía con sólida formación en análisis cuantitativo y ciencia de datos. Apasionado por descubrir patrones en datos y convertir números en historias que impulsen decisiones. Con experiencia académica y proyectos personales en análisis de datos económicos, visualización y programación estadística.

	%Competente en procesamiento y limpieza de bases de datos, análisis exploratorio, visualización con ggplot2 y Matplotlib, y elaboración de reportes automatizados. Familiarizado con SQL para consultas de bases de datos y nociones de Machine Learning aplicado. Rápido aprendizaje de nuevas tecnologías y frameworks.

	%Busco mi primera oportunidad profesional en análisis de datos donde pueda aplicar mis conocimientos técnicos, continuar desarrollándome en un entorno colaborativo y contribuir con iniciativa, dedicación y capacidad analítica al éxito de proyectos basados en datos.

	%%%%%%%%%%%%%%%%%%%%%%%%%%%%%%%%%%%%%%%%
	% VERSIÓN ALTERNATIVA - Data Engineer
	%%%%%%%%%%%%%%%%%%%%%%%%%%%%%%%%%%%%%%%%
	%Economista con especialización en ingeniería de datos y automatización de procesos analíticos. Experiencia en diseño de pipelines ETL, gestión de bases de datos relacionales (MySQL, PostgreSQL) y optimización de flujos de trabajo con datos estructurados y no estructurados.

	%Dominio de Python para automatización, R para análisis estadístico y SQL para manipulación eficiente de grandes volúmenes de datos. Familiarizado con principios de arquitectura de datos, versionamiento de código (Git) y buenas prácticas de documentación técnica.

	%Busco oportunidades en Data Engineering o roles técnicos de datos donde pueda diseñar infraestructuras escalables, optimizar procesos de datos y colaborar con equipos multidisciplinarios para construir soluciones robustas que soporten decisiones basadas en evidencia.

	%%%%%%%%%%%%%%%%%%%%%%%%%%%%%%%%%%%%%%%%
	% VERSIÓN ALTERNATIVA - Research Data Analyst
	%%%%%%%%%%%%%%%%%%%%%%%%%%%%%%%%%%%%%%%%
	%Economista con formación metodológica rigurosa y especialización en análisis de datos para investigación. Experto en diseño de estudios cuantitativos, análisis estadístico avanzado (regresiones, series de tiempo, análisis multivariado) y visualización de resultados de investigación con estándares académicos.

	%Dominio de software estadístico (R, Stata, SPSS) para procesamiento de encuestas, análisis de datos panel, evaluación de impacto y estudios observacionales. Experiencia en redacción de reportes técnicos, papers académicos y presentación de hallazgos a audiencias especializadas.

	%Busco oportunidades en centros de investigación, think tanks o unidades de estudios donde pueda aplicar mi formación en economía y estadística para desarrollar investigaciones aplicadas, evaluar políticas públicas y contribuir a la generación de evidencia para la toma de decisiones informadas.

}
